% ============================================================================
\section{Introduction}
\label{sec:intro}

Looped language models reuse the same block across recurrent depth.\footnote{Code: \url{https://github.com/Rituraj003/readout-blind-spot}.}
This makes test-time depth a natural compute knob: stop after a few loops when the prediction is easy, or continue iterating when additional computation helps.
Universal Transformers \citep{dehghani2019universal}, Huginn \citep{geiping2025huginn}, Ouro \citep{zhu2025ouro}, LoopFormer \citep{jeddi2026loopformer}, and Elastic Looped Transformer \citep{goyal2026elt} all explore versions of this idea.

Looping changes the role of a hidden state.
In a fixed-depth transformer, a late hidden state is mostly an interface to the output head.
In a looped model, the same hidden state is also runtime state: it is decoded for prediction and then reused as input to future computation.
This makes a simple question unavoidable:
\emph{which parts of the recurrent state does the supervised loss actually control?}

Dense supervision is the natural answer.
Apply cross-entropy (CE) at every loop, and each intermediate state receives direct prediction supervision.
This indeed trains the prediction interfaces: early exits become usable.
But dense CE does not automatically control every variable carried by the recurrent state.
It controls what the readout makes visible.
If an active recurrent variable is hidden by every supervised readout, the model can be trained at every loop while that variable remains underconstrained.

We study this readout blind spot for hidden-state scale.
In pre-norm residual loops, scale is \emph{active}: the residual skip carries it forward and the learned update can change it.
With an RMSNorm or LayerNorm readout, scale is \emph{locally invisible}: multiplying the hidden state by a positive scalar leaves the readout unchanged up to the normalizer epsilon, so the immediate CE loss has essentially no radial derivative.
The result is a visibility--activity mismatch.
In our looped transformers, per-loop CE through normalized readouts still allows hidden-state norms to drift by orders of magnitude when scale remains active in the recurrent path.
The model is supervised at every loop, yet every supervised interface hides the variable that is drifting.

This yields the central design rule of the paper:
\begin{quote}
    \emph{Per-loop cross-entropy trains exit usability; recurrent scale control requires either making hidden-state scale visible to a loss or removing it from the recurrent path.}
\end{quote}

\begin{figure}[t]
\centering
\includegraphics[width=\linewidth]{figures/endpoint_scale_quadrant.pdf}
\caption{\textbf{Exit training and scale control are orthogonal in the core 2$\times$2.}
The x-axis measures early-exit failure, the cross-entropy gap between the first and fourth recurrent loops, $\mathrm{CE}(K{=}1)-\mathrm{CE}(K{=}4)$, where $K$ is the inference loop count; the axis uses a symlog scale.
The y-axis measures final-loop norm drift, with $H_K$ denoting the hidden state after $K$ loops.
Lower is better on both axes.
Per-loop CE moves models left by making intermediate exits usable, while scale-visible readouts move models down by controlling recurrent scale.
The desired region requires both interventions; later norm-penalty and final-only-normalization controls show that the effect is not merely a raw-readout artifact.}
\label{fig:decomp}
\end{figure}

We support this rule with a compact mechanism and controlled ablations.
The mechanism has two components.
Scale-invariant readouts such as RMSNorm or LayerNorm produce zero immediate radial cross-entropy signal, while pre-norm residual recurrence continues to preserve and update scale.
Future-loop losses can still carry indirect radial information, but that signal weakens at large hidden-state norms.
The experiments then test the interventions suggested by this mechanism: exposing scale through raw readouts or explicit norm penalties, or removing scale through normalization inside the recurrent path.

Our contributions are:
\begin{enumerate}
    \item We introduce the readout blind spot as a visibility--activity mismatch: dense CE controls variables exposed by the readout, not every variable active in recurrence.
    In 44M- and 129M-parameter looped language models where hidden-state scale remains active in the recurrent path, per-loop cross-entropy through normalized readouts still produces large norm drift.
    
    \item We instantiate the blind spot for hidden-state scale. Normalized readouts hide scale from the immediate cross-entropy loss, while pre-norm residual recurrence continues to propagate it. Raw readouts, norm penalties, final-only normalization, gradient diagnostics, and recurrent-scale clamps all support this interpretation.
    
    \item We show that scale control improves the variable-depth perplexity--compute frontier. Scale-controlled per-loop models maintain usable exits while achieving lower perplexity than the $K$-invariant RMSNorm baseline at matched inference-depth operating points in our benchmarks.
\end{enumerate}

